\documentclass[11pt,a4paper]{article}
\usepackage[T1]{fontenc}
\usepackage[utf8]{inputenc}
\usepackage[french]{babel}
\usepackage{lmodern}
\usepackage[margin=2.5cm]{geometry}
\usepackage{amsmath,amssymb}
\usepackage{hyperref}

% --- Articles et correspondance avec la formalisation ----------------------
% \article[niveau]{Titre}    niveau d'impérativité (M4) : absolue, relative
%                            ou dispositive ; optionnel (article sans norme,
%                            par exemple une définition).
% \lean{Nom}                 déclaration Lean qui formalise l'article.
% \theoreme{Nom}             théorème Lean qui prouve une propriété de l'article.
% \horsformalisation{Raison} article volontairement non formalisé.
% Les noms Lean s'écrivent sans le préfixe « Prouvhommes. ».

\newcounter{article}
\newcommand{\niveau}[1]{\hfill{\footnotesize\textsc{#1}}}
\newcommand{\article}[2][]{%
  \refstepcounter{article}%
  \par\medskip\noindent\textbf{Article \thearticle\ --- #2}%
  \if\relax\detokenize{#1}\relax\else\niveau{#1}\fi
  \par\nopagebreak}

\newcommand{\renvoi}[2]{%
  \par\nopagebreak\noindent{\footnotesize\textsf{#1 :} \texttt{\detokenize{#2}}}\par}
\newcommand{\lean}[1]{\renvoi{Lean}{#1}}
\newcommand{\theoreme}[1]{\renvoi{Théorème}{#1}}
\newcommand{\horsformalisation}[1]{%
  \par\nopagebreak\noindent{\footnotesize\textsf{Hors formalisation :} #1}\par}

\title{Prouv'hommes\\[0.3em]\large Proposition de code du travail simplifié}
\author{}
\date{Version de travail --- \today}

\begin{document}
\maketitle
\tableofcontents

\section*{Préambule}
Ce code tient en une cinquantaine d'articles. Chaque article a une
contrepartie formelle dans le dossier \texttt{lean/} du dépôt.

\section{Titre I --- Le contrat de travail}

\article[absolue]{Forme unique}\label{art:forme-unique}
Le contrat de travail a une seule forme. Il est conclu sans détermination
de durée. Toute clause fixant un terme est réputée non écrite.
\lean{Contrat}
\theoreme{pas_de_terme}

\article[relative]{Écrit}\label{art:ecrit}
Le contrat est constaté par écrit. L'écrit indique les parties, la date de
début, la fonction, le salaire mensuel et la durée de travail convenue.
\lean{Contrat}

\article[dispositive]{Période d'essai}\label{art:essai}
Les trente premiers jours du contrat constituent une période d'essai. Un
accord d'entreprise ou le contrat peut la réduire, la supprimer ou l'allonger,
dans la limite absolue de quatre-vingt-dix jours.
\lean{Contrat.dureeEssai}
\lean{dispDureeEssai}
\theoreme{dureeEssai_le_max}


\end{document}
